%***************************************************************************************************
% preamble/macros.tex - makra pomocnicze, środowiska matematyczne, pytania i odpowiedzi, listy.
%***************************************************************************************************

% ---------------------------------------------------------------------------------------------
% Makra semantyczne
% ---------------------------------------------------------------------------------------------
\newcommand{\code}[1]{\texttt{\detokenize{#1}}}
\newcommand{\file}[1]{\texttt{\detokenize{#1}}}
\newcommand{\cmd}[1]{\texttt{\detokenize{#1}}}
\newcommand{\latexcmd}[1]{\texttt{\textbackslash #1}}
\newcommand{\hash}[1]{{\ttfamily\footnotesize\seqsplit{#1}}}

\newcommand{\lead}[1]{\par\medskip\noindent\textbf{#1}\par\nopagebreak\smallskip}

\newcommand{\UnnumberedChapter}[1]{%
  \clearpage\phantomsection
  \addcontentsline{toc}{chapter}{#1}%
  \chapter*{#1}%
}

\newcounter{qacount}
\newenvironment{qa}[1]{%
  \refstepcounter{qacount}%
  \par\medskip\noindent\textbf{Pytanie \theqacount:} \textit{#1}\par\nopagebreak\smallskip
  \noindent\textbf{Odpowiedź}\par\nopagebreak
}{\par\medskip}

\newenvironment{abbreviations}
  {\begin{description}[style=multiline, leftmargin=3.5cm, font=\bfseries, itemsep=2pt]}
  {\end{description}}

\newcommand{\textapprox}{\raisebox{0.5ex}{\texttildelow}}

\newcommand{\cmark}{\ding{51}}
\newcommand{\xmark}{\ding{55}}
\newcommand{\hcell}[1]{\textbf{#1}}

% ---------------------------------------------------------------------------------------------
% Skróty matematyczne
% ---------------------------------------------------------------------------------------------
\newcommand{\NN}{\mathbb{N}}
\newcommand{\ZZ}{\mathbb{Z}}
\newcommand{\QQ}{\mathbb{Q}}
\newcommand{\RR}{\mathbb{R}}
\newcommand{\CC}{\mathbb{C}}
\newcommand{\IMP}{\rightarrow}
\newcommand{\IFF}{\leftrightarrow}

% ---------------------------------------------------------------------------------------------
% Twierdzenia, definicje, przykłady, uwagi
% ---------------------------------------------------------------------------------------------
\theoremstyle{plain}
\newtheorem{theorem}{Twierdzenie}[chapter]
\newtheorem{lemma}[theorem]{Lemat}
\newtheorem{corollary}[theorem]{Wniosek}

\theoremstyle{definition}
\newtheorem{definition}[theorem]{Definicja}
\newtheorem{example}[theorem]{Przykład}

\theoremstyle{remark}
\newtheorem{remark}[theorem]{Uwaga}

\renewcommand{\qedsymbol}{$\blacklozenge$}

% ---------------------------------------------------------------------------------------------
% Listy
% ---------------------------------------------------------------------------------------------
\setlist{topsep=0.4em, itemsep=0.2em, parsep=0pt, partopsep=0pt}
\setlist[itemize,1]{label=\textbullet, leftmargin=*}
\setlist[itemize,2]{label=-, leftmargin=*}
\setlist[itemize,3]{label=$\diamond$, leftmargin=*}
\setlist[enumerate,1]{label=\arabic*., leftmargin=*}
\setlist[enumerate,2]{label=\alph*), leftmargin=*}

\setlength{\emergencystretch}{2em}
